\chapter{Discusión}
\label{ch:discusion}

Este capítulo interpreta la trayectoria experimental completa presentada en el Capítulo~\ref{ch:resultados}. La tesis no culmina en una comparación aislada TRSS--MNE: antes de esa comparación se implementó un espacio amplio de métodos, se realizó una selección preliminar de combinaciones método--grafo, se redujeron candidatos por criterios metodológicos y computacionales, y se congelaron variantes para evitar que la comparación final dependiera de decisiones posteriores al resultado. La discusión se organiza en cuatro niveles: el valor científico de la trayectoria experimental, la interpretación condicional de TRSS frente a MNE, las tensiones entre métricas y costo, y las limitaciones que deben quedar explícitas antes de formular conclusiones.

\section{Lectura integrada de la trayectoria experimental}

El primer resultado conceptual del trabajo es que la evaluación de interpolación EEG no puede reducirse a escoger el método con menor error promedio. La fase inicial mostró un espacio heterogéneo: algunas familias fueron competitivas en datos sintéticos, otras en registros EEG reales, y varias dependieron fuertemente de la construcción del grafo, de la disponibilidad de coordenadas o del costo de ejecución. Por eso, la reducción de candidatos no debe entenderse como una etapa administrativa, sino como parte del resultado experimental: documenta qué opciones eran plausibles, cuáles se descartaron y por qué la comparación final se concentra en TRSS y MNE-Python.

Esta separación evita confundir evidencia preliminar con evidencia final. Los rankings exploratorios de la Tabla~\ref{tab:ch4_preliminary_selection_top_candidates} orientaron el diseño, pero no se usaron como conclusión principal. La evidencia final proviene de comparaciones pareadas sobre la misma señal, máscara y semilla, resumidas en las tablas y figuras posteriores del Capítulo~\ref{ch:resultados}. Así, la tesis no selecciona TRSS de forma arbitraria: llega a TRSS mediante una secuencia trazable de implementación, selección preliminar, reducción, optimización y comparación final.

\section{Aporte de la selección preliminar de métodos y grafos}

La selección preliminar aportó tres lecciones. Primero, las variantes con información temporal aparecieron de forma recurrente entre las combinaciones competitivas. Esto justificó estudiar con mayor detalle una formulación espacio--temporal, en vez de limitar la tesis a interpolación puramente espacial. Segundo, los grafos geométricos reproducibles resultaron preferibles para sostener conclusiones finales, porque dependen de posiciones de electrodos y no de la misma señal que se intenta reconstruir. Tercero, algunos grafos aprendidos o basados en correlación fueron informativos como exploración, pero no adecuados como soporte principal de la comparación final por costo, reproducibilidad o riesgo de circularidad.

Esta lectura rescata el valor de los resultados intermedios sin exagerar su alcance. Un método descartado no es necesariamente inútil; puede haber sido demasiado costoso, difícil de estabilizar, sensible a hiperparámetros o poco apropiado para una comparación final controlada. En ese sentido, la fase de reducción cumple una función científica: transforma un espacio experimental amplio en una comparación interpretable, auditable y menos propensa a sobreajuste narrativo.

También es importante que el documento mantenga visibles los métodos no finalistas. Tikhonov, BGSRP, TV en grafos, ICA/FastICA/Picard, RBF, IDW, vecino más cercano, interpolación lineal y splines esféricos forman parte del espacio experimental que condujo a la comparación final. Que no todos aparezcan con igual profundidad en el cierre no significa que hayan sido ignorados, sino que la tesis distingue entre cobertura exploratoria y evidencia final.

\section{Por qué MNE-Python sigue siendo una referencia fuerte}

MNE-Python debe interpretarse como una referencia comunitaria madura, no como una línea base débil. Su interpolación por splines esféricas combina un prior espacial fisiológicamente razonable con una implementación robusta, normalización interna, regularización numérica y uso global de los canales disponibles. Cuando la pérdida es baja, el vecindario espacial conserva suficiente información o la señal es espacialmente suave, este prior puede reconstruir canales faltantes sin requerir una restricción temporal adicional.

Esta fortaleza explica los empates y los casos favorables a MNE observados en el Capítulo~\ref{ch:resultados}. También eleva el estándar de comparación: una ventaja de TRSS frente a MNE ocurre contra una herramienta real usada en la práctica, no contra una implementación simplificada diseñada solo para la tesis. Por ello, MNE debe mantenerse como primera referencia para estudios EEG y como opción operativa cuando rapidez, estabilidad y reproducibilidad estándar sean criterios dominantes.

\section{Interpretación condicional de TRSS}

TRSS aporta valor cuando la información espacial disponible se vuelve insuficiente. En pérdidas cercanas, periféricas o severas, el canal oculto puede quedar sin vecinos informativos; en ese régimen, la continuidad temporal restringe el conjunto de soluciones plausibles y puede reducir errores de amplitud o forma. Esta es la zona donde los resultados descriptivos muestran la utilidad más clara de TRSS.

Sin embargo, la conclusión no es una dominancia universal. En escenarios de baja pérdida, señales suaves o restricciones fuertes de tiempo, MNE puede igualar o superar a TRSS. Además, el ajuste de hiperparámetros condiciona el desempeño de la regularización temporal. La contribución debe formularse, por tanto, como una frontera práctica de uso: TRSS es justificable cuando la pérdida espacial amenaza la calidad de reconstrucción y el análisis admite costo adicional; MNE sigue siendo preferible en condiciones estándar o cuando la velocidad pesa más que una ganancia moderada de reconstrucción.

\section{Tensión entre precisión temporal, espectro y costo}

Las métricas del Capítulo~\ref{ch:resultados} miden propiedades distintas. MAE, RMSE y NRMSE resumen error de amplitud; DTW, correlación y $R^2$ se relacionan con forma temporal; LSD, coherencia y PSD capturan preservación espectral; el tiempo de ejecución determina factibilidad práctica. Por eso, mejorar una familia de métricas no implica mejorar todas. TRSS puede acercarse mejor punto a punto a la señal original y aun así no dominar la distancia espectral global.

Esta tensión obliga a elegir el método según el objetivo posterior. Si la prioridad es reconstruir segmentos para análisis temporal fuera de línea, TRSS puede ser atractivo en escenarios difíciles. Si la prioridad es preservar propiedades espectrales globales, procesar grandes volúmenes de datos o mantener un flujo estándar de preprocesamiento, MNE conserva ventajas claras. Las figuras temporales y espectrales del Capítulo~\ref{ch:resultados} no sustituyen una validación por tarea, pero sí muestran por qué una única métrica sería insuficiente para decidir.

\section{Costo computacional y escalabilidad}

El costo computacional es una limitación práctica de TRSS. Aunque los tiempos por caso pueden ser aceptables en análisis planificados y fuera de línea, la diferencia frente a MNE se vuelve relevante en estudios con muchos sujetos, ventanas temporales, variantes de grafo o búsqueda de hiperparámetros. En aplicaciones interactivas, pipelines clínicos o procesamiento masivo, esta diferencia puede pesar más que la mejora en una métrica de reconstrucción.

La escalabilidad futura requiere cambios de ingeniería: matrices dispersas, factorizaciones reutilizables del Laplaciano, solvers iterativos con parada temprana, paralelización por sujeto o por bloque de canales, y reglas de hiperparámetros que eviten búsqueda exhaustiva. Sin estas mejoras, TRSS debe entenderse como una alternativa de análisis controlado, no como sustituto inmediato de MNE en flujos rutinarios de baja latencia.

\section{Deficiencias y carencias del estudio}

La principal limitación experimental es el uso de canales ocultos artificialmente. Esta estrategia permite disponer de verdad de terreno y medir error de reconstrucción, pero no reproduce todos los casos reales de canal dañado: saturación, ruido muscular, deriva lenta, contactos inestables o artefactos intermitentes pueden afectar de manera distinta a cada método. Por ello, los resultados son evidencia controlada de interpolación bajo ausencia simulada, no validación clínica completa.

Una segunda carencia es la ausencia de evaluación experta ciega. La plausibilidad visual de las reconstrucciones se inspecciona mediante señales representativas y PSD, pero no se sometió a revisión sistemática por especialistas EEG. Tampoco se evaluó el impacto en tareas posteriores como clasificación BCI, detección de potenciales evocados, conectividad, análisis por bandas o localización de fuentes. Un método que reduce MAE no necesariamente mejora una decisión neurofisiológica posterior.

También debe reconocerse que la reducción del espacio experimental dejó familias fuera de la comparación final. Los grafos aprendidos, grafos adaptativos y métodos temporales alternativos quedaron como evidencia intermedia o trabajo futuro, no como conclusiones cerradas. Esta decisión fue necesaria para evitar una comparación saturada y difícil de interpretar, pero implica que la tesis no agota todas las combinaciones posibles de interpolación y construcción de grafos.

Finalmente, la selección de hiperparámetros sigue siendo una limitación abierta. El oráculo representa una cota superior no desplegable; la variante calibrada reduce parcialmente ese problema, pero todavía se requiere una regla automática basada en propiedades observables del registro, la máscara de canales faltantes, el grafo y la señal. Sin esa regla, una parte del potencial de TRSS depende de calibraciones costosas.

\section{Aspectos fuera del alcance}

Algunas decisiones quedaron fuera por tiempo, recursos o alcance metodológico. No se incorporaron modelos supervisados ni aprendizaje profundo porque habrían requerido particiones de entrenamiento, validación externa y controles adicionales para evitar fuga de información. No se hizo evaluación clínica con expertos ni validación sobre canales realmente dañados, porque ello requiere anotaciones externas y criterios de calidad distintos a la verdad de terreno simulada. Tampoco se evaluó desempeño en tareas posteriores, ya que cada tarea introduce métricas, pipelines y fuentes de variabilidad propias.

Del mismo modo, no se agotaron grafos aprendidos ni adaptativos. Su inclusión rigurosa exigiría separar construcción del grafo y evaluación de reconstrucción para evitar circularidad, además de medir estabilidad entre ejecuciones y costo. La tesis optó por concentrar la interpretación principal en una comparación más acotada y reproducible. Esta decisión reduce amplitud, pero aumenta interpretabilidad.

\section{Lecciones metodológicas}

La primera lección es que una referencia comunitaria fuerte debe evaluarse con el mismo cuidado que el método propuesto. Comparar TRSS solo contra líneas base internas habría producido una lectura incompleta. MNE-Python obliga a reconocer los escenarios donde una herramienta madura sigue siendo superior.

La segunda lección es que los descartes también son resultados. Mostrar la selección preliminar y la reducción de candidatos hace más transparente el proceso experimental y evita que el lector vea la comparación final como una decisión posterior no justificada. En trabajos computacionales con muchas variantes, explicar qué no se usó y por qué es parte de la contribución.

La tercera lección es que la evaluación multi-métrica es indispensable. Amplitud, forma temporal, espectro y costo no colapsan en una única respuesta. La tesis gana solidez porque explicita ganancias, pérdidas y empates prácticos, y porque conecta cada afirmación con artefactos reproducibles: tablas, CSV, figuras, scripts y validadores.

En síntesis, la discusión no transforma los resultados en una recomendación universal, sino en una regla de lectura. MNE-Python debe seguir siendo la referencia operacional cuando el escenario es favorable a la suavidad espacial o cuando el costo domina. TRSS se vuelve relevante cuando la pérdida espacial deteriora la vecindad local y la continuidad temporal aporta una restricción adicional plausible. Esta distinción es el puente natural hacia las conclusiones del capítulo siguiente.
